% !PRINTED 635
\weberpage{635}
%<<CONT
\weq{5}{\begin{aligned}
(\varepsilon^{g}, x)(\varepsilon, x)^{-g} &= f_0\\
(\varepsilon^{g^2}, x)(\varepsilon^{g}, x)^{-g} &= f_1\\
\cdot\ \cdot\ \cdot\ \cdot\ \cdot\ \cdot\ \cdot\ \cdot\ \cdot\ \cdot & \\
(\varepsilon^{g^{n-1}}, x)(\varepsilon^{g^{n-2}}, x)^{-g} &= f_{n-2}.
\end{aligned}}
so dass also, wenn $(\varepsilon^{g}, x)(\varepsilon, x)^{-g} = f(\varepsilon)$ gesetzt wird,
\weq{6}{f_h = f(\varepsilon^{g^h})}
ist. Wendet man auf (5) die Permutation $t^{-1}$ an, so erleiden, wie nach dem Theorem 4. zu sehen ist, die Functionen
\[
f_0, f_1, f_2 \ldots f_{n-2}
\]
eine cyklische Permutation $(0, 1, 2 \ldots n-2)$.

Ausserdem sind die Functionen $f_\nu$, so lange die $x$ variabel sind, wie man aus (5) ersieht, alle \gesperrt{von einander verschieden} da sie in verschiedene lineare Factoren zerlegt sind.

Bilden wir also eine cyklische Function dieser $n-1$ Grössen
\weq{7}{C(f_0, f_1, f_2 \ldots f_{n-2}),}
so bleibt diese sowohl durch $s$ als durch $t^{-1}$, also durch $t$ und die ganze metacyklische Gruppe ungeändert, und wir erhalten das Theorem:

\gesperrt{6. Stellt man eine cyklische Function der Grössen $f_0, f_1 \ldots f_{n-1}$ in der Normalform dar, so sind die Coëfficienten metacyklische Functionen der Variablen $x_0, x_1 \ldots x_{n-1}$.} % [sic] Print says $f_{n-1}$ although (5) defines the sequence only through $f_{n-2}$; preserved.

Hierbei ist unter einer metacyklischen Function jede Function zu verstehen, die durch die Permutationen der linearen Gruppe ungeändert bleibt.

Die Functionen $f_0, f_1 \ldots f_{n-2}$ gehen aber auch cyklisch in einander über, wenn man die $x$ ungeändert lässt, und für $\varepsilon$ die Substitution $(\varepsilon, \varepsilon^g)$ macht. Wenn also die cyklische Function $C$ von den $f$ in ihren Coëfficienten $\varepsilon$ nicht enthält, so bleibt $C$ durch sämmtliche Substitutionen $(\varepsilon, \varepsilon^h)$ ungeändert, und wir erhalten aus 3. das Theorem:

\gesperrt{7. Ist $C(f_0, f_1 \ldots f_{n-2})$ eine cyklische Function der Grössen $f_0, f_1 \ldots f_{n-2}$ mit rationalen Coëfficienten, so ist $C$ eine metacyklische Function der Variablen $x_0, x_1 \ldots x_{n-1}$ mit rationalen Coëfficienten.}

Wenn wir die Functionen $f_0, f_1 \ldots f_{n-2}$, wie sie durch (5) gegeben sind, der Reihe nach zu den Potenzen $g^{n-2}, g^{n-3} \ldots 1$
%CONT>>
